\section*{Capítulo \ref{ch:b2:mammal-reproduction} --- Reproducción sexual de los mamíferos}

\begin{solution}{exo:b2:mammal-reproduction:1}
Las espermatogonias se dividen desde la pubertad y de por vida; las
ovogonias, solo en el feto. Una \omterm{def:b2:meiosis-heredity:meiosis}{meiosis} da cuatro espermatozoides, pero un
solo óvulo (y tres cuerpos polares). Un espermatozoide tarda unos 64 días
más dos semanas de maduración; un óvulo queda detenido en \omterm{prop:b2:meiosis-heredity:stages}{profase I} desde
antes del nacimiento hasta el mes de su \omterm{def:b2:mammal-reproduction:ovary}{ovulación} --- décadas --- y de
nuevo en metafase II hasta la fecundación. Un hombre produce $10^{8}$
espermatozoides al día, unos $10^{12}$ a lo largo de su vida; una mujer
ovula unos 400 óvulos.
\end{solution}

\begin{solution}{exo:b2:mammal-reproduction:2}
Unión a la \omterm{prop:b2:mammal-reproduction:fertilisation}{zona pelúcida}; \omterm{prop:b2:mammal-reproduction:fertilisation}{reacción acrosómica} y digestión de un camino;
fusión de las membranas y entrada del núcleo del espermatozoide; onda de
calcio y \omterm{prop:b2:mammal-reproduction:fertilisation}{reacción cortical} (endurecimiento de la zona y pérdida de sus
receptores: el bloqueo de la poliespermia); finalización de la \omterm{def:b2:meiosis-heredity:meiosis}{meiosis} II
y expulsión del segundo cuerpo polar; formación de los dos pronúcleos,
replicación del ADN, primera mitosis.
\end{solution}

\begin{solution}{exo:b2:mammal-reproduction:3}
El \omterm{def:b2:mammal-reproduction:implantation}{trofoblasto} (capa externa) forma la \omterm{def:b2:mammal-reproduction:placenta}{placenta} y las membranas; la \omterm{def:b2:mammal-reproduction:implantation}{masa
celular interna} forma el embrión. La \omterm{def:b2:mammal-reproduction:implantation}{implantación} empieza los días 6 a 7
tras la fecundación, en un revestimiento uterino preparado por la
progesterona del \omterm{def:b2:mammal-reproduction:ovary}{cuerpo lúteo}.
\end{solution}

\begin{solution}{exo:b2:mammal-reproduction:4}
Oxígeno: difusión. Glucosa: transporte facilitado. Glóbulos rojos
maternos: no atraviesan (las sangres nunca se mezclan). IgG: transporte
mediado por receptores. Urea: difusión, del feto a la madre. Alcohol:
difunde libremente. Bacterias: la mayoría no atraviesan (algunas, como la
de la sífilis o la listeria, sí).
\end{solution}

\begin{solution}{exo:b2:mammal-reproduction:5}
$m = 3, 1, 0.4, 0.1$: $P = 0.95, 0.63, 0.33, 0.095$. Por debajo de un
medio cuando $m < \ln 2 = 0.69$, es decir, por debajo de $7\times 10^{7}$
espermatozoides. La curva se satura: por encima de unos $10^{8}$, los
espermatozoides adicionales apenas añaden nada; por debajo, la fertilidad
cae casi en proporción al recuento.
\end{solution}

\begin{solution}{exo:b2:mammal-reproduction:6}
Antes de la pubertad: $7\times 10^{5}$ perdidos en 13 años (\num{4750}
días): unos 150 al día. Después: \num{299000} perdidos en 37 años
(\num{13500} días): 22 al día; \omterm{def:b2:mammal-reproduction:ovary}{ovulación}: $400/\num{13500} = 0.03$ al día.
La atresia supera a la \omterm{def:b2:mammal-reproduction:ovary}{ovulación} en una proporción de setecientos a uno;
la reserva no se agota ovulando, sino muriendo.
\end{solution}

\begin{solution}{exo:b2:mammal-reproduction:7}
Adulta: $(30/27)^{2.7} = 1.33$, $S = 1.33/2.33 = 0.57$. Fetal:
$(30/19)^{2.7} = 3.4$, $S = 3.4/4.4 = 0.77$. A las bajas presiones de la
\omterm{def:b2:mammal-reproduction:placenta}{placenta}, la sangre fetal carga una quinta parte más de oxígeno que la
sangre adulta, y lo descarga en los tejidos a presiones en las que la
hemoglobina adulta ya estaría medio vacía.
\end{solution}

\begin{solution}{exo:b2:mammal-reproduction:8}
$2^{k} = 25$: $k = 4.6$ duplicaciones, \qty{9.3}{d}: días 16 a 17 tras la
fecundación; días 30 a 31 tras la última regla --- aproximadamente el día
en que falta la regla.
\end{solution}

\begin{solution}{exo:b2:mammal-reproduction:9}
$3\times 10^{-8}\times 6\times 10^{4}\times 20/10^{-3} =
\qty{36}{mL/min}$ frente a una demanda de \qty{25}{mL/min}: el margen pasa
de ocho a uno y medio; el feto sufre hipoxia ante cualquier carga
adicional, crece despacio y puede requerir un parto anticipado.
\end{solution}

\begin{solution}{exo:b2:mammal-reproduction:10}
(a) Se forman \omterm{def:b2:mammal-reproduction:testis}{testículos} que secretan ambas hormonas; la AMH elimina los
conductos de Müller, pero la testosterona no puede actuar: los conductos
de Wolff regresan y los genitales externos son femeninos --- una mujer
con \omterm{def:b2:mammal-reproduction:testis}{testículos} y sin útero. (b) La testosterona actúa y la AMH no: un
varón con epidídimo, conducto deferente y genitales masculinos, y además
útero y oviductos. (c) Ovarios y derivados de Müller (sin \omterm{def:b2:mammal-reproduction:testis}{testículo}, sin
AMH); los andrógenos suprarrenales virilizan los genitales externos en
grado variable. (d)
\emph{SRY} forma \omterm{def:b2:mammal-reproduction:testis}{testículos}, que imponen la vía masculina: un hombre con
dos cromosomas X, estéril porque faltan los genes del Y necesarios para
la \omterm{def:b2:mammal-reproduction:testis}{espermatogénesis}.
\end{solution}

\begin{solution}{exo:b2:mammal-reproduction:11}
La más frecuente (\qty{68}{\%}) es la división entre los días 4 y 8: tras
formarse el \omterm{def:b2:mammal-reproduction:implantation}{trofoblasto} (de ahí una \omterm{def:b2:mammal-reproduction:placenta}{placenta}), pero antes del amnios (de
ahí dos sacos). Antes del día 3 cada mitad forma su propio \omterm{def:b2:mammal-reproduction:implantation}{trofoblasto}, y
de ahí dos \omterm{def:b2:mammal-reproduction:placenta}{placentas} (\qty{30}{\%}); después del día 8 el amnios ya está
formado y se comparte (\qty{2}{\%}). Dividir la \omterm{def:b2:mammal-reproduction:implantation}{masa celular interna} una
vez que el \omterm{def:b2:mammal-reproduction:implantation}{trofoblasto} está determinado es, evidentemente, el suceso más
fácil.
\end{solution}

\begin{solution}{exo:b2:mammal-reproduction:12}
Un marsupial invierte poco antes del nacimiento: una gestación corta, una
cría diminuta y una lactancia que puede interrumpirse desprendiéndose de
la cría de la bolsa si llega la sequía, de modo que el riesgo de la madre
en cada intento es bajo y puede volver a intentarlo enseguida. Un
placentario invierte mucho en una gestación larga, con una cría grande y
bien desarrollada que sobrevive mejor al nacer, a costa de una madre
comprometida durante meses y vulnerable mientras tanto. En ambientes
productivos y previsibles gana la mayor supervivencia de las crías
placentarias; el clima errático de Australia premia la opción de
abandonar de los marsupiales, y su aislamiento mantuvo fuera a los
placentarios hasta hace poco.
\end{solution}

\begin{solution}{pb:b2:mammal-reproduction:1}
\textbf{1.} $m = 3\times 10^{8}\times 10^{-8} = 3$; $P = 1 - e^{-3} =
0.95$.
\textbf{2.} $1 - e^{-3}(1 + 3) = 0.80$: sin el bloqueo, cuatro de cada
cinco fecundaciones serían poliespérmicas y el embrión triploide.
\textbf{3.} $0.15/5\times 10^{-5} = \qty{3000}{s}$, cincuenta minutos;
los transportan las contracciones del útero y del oviducto.
\textbf{4.} $300/3\times 10^{8} = 10^{-6}$: se pierden en la vagina
ácida, en el moco cervical, a manos de los fagocitos del útero y en el
oviducto equivocado.
\textbf{5.} Desde tres días antes de la \omterm{def:b2:mammal-reproduction:ovary}{ovulación} hasta un día después:
unos cuatro días.
\textbf{6.} $m = 0.5$, $P = 0.39$; ciclos esperados $1/P = 2.5$.
\textbf{7.} Las pocas mitocondrias del espermatozoide, situadas en la
pieza intermedia, se marcan y se destruyen tras su entrada; el óvulo
lleva cien mil.
\textbf{8.} $120/16 = 7.5$: siete divisiones, $2^{7} = 128$ células ---
alrededor de un centenar.
\textbf{9.} $\log_2 25 = 4.6$ duplicaciones, \qty{9.3}{d}: días 16 a 17;
días 30 a 31 tras la regla.
\textbf{10.} El \omterm{def:b2:mammal-reproduction:ovary}{cuerpo lúteo} regresa unos doce días después de la
\omterm{def:b2:mammal-reproduction:ovary}{ovulación} salvo que la hCG lo rescate; sin ella la progesterona cae, el
revestimiento se desprende y el embrión se pierde con la regla.
\textbf{11.} $\log_2 10^{5} = 16.6$ duplicaciones; hacia la semana 10 la
\omterm{def:b2:mammal-reproduction:placenta}{placenta} produce su propia progesterona y el \omterm{def:b2:mammal-reproduction:ovary}{cuerpo lúteo} deja de ser
necesario.
\textbf{12.} Los óvulos permanecen décadas en \omterm{prop:b2:meiosis-heredity:stages}{profase I} y las cohesinas
que mantienen unidos los bivalentes se degradan, de modo que la no
disyunción en la \omterm{def:b2:meiosis-heredity:meiosis}{meiosis} I aumenta bruscamente con la edad materna; los
espermatozoides se fabrican de nuevo en dos meses.
\textbf{13.} Tras el \omterm{def:b2:mammal-reproduction:implantation}{trofoblasto} (día 5) y antes del amnios (día
8): una \omterm{def:b2:mammal-reproduction:placenta}{placenta} y dos sacos.
\textbf{14.} $3\times 10^{-8}\times 1.2\times 10^{5}\times 20/3.5\times
10^{-4} = \qty{206}{mL/min}$.
\textbf{15.} $7\times 3.5 = \qty{24.5}{mL/min}$: un margen de ocho.
\textbf{16.} $500\times 0.16 = \qty{80}{mL/min}$; extracción de
$24.5/80 = \qty{31}{\%}$.
\textbf{17.} $5\times 3.5\times 1440 = \qty{25}{g/d}$; contenidos en
\qty{28}{L} de sangre materna; flujo diario de \qty{720}{L}: el feto toma
alrededor del \qty{4}{\%} de la glucosa que pasa.
\textbf{18.} La hemoglobina fetal está saturada al \qty{77}{\%} a
\qty{30}{mmHg}, el feto tiene más glóbulos rojos y un gasto cardíaco
elevado, y sus tejidos están adaptados a trabajar a baja presión --- vive,
en la práctica, en la cima del Everest.
\textbf{19.} El intercambio de gases sin un pulmón funcional, la nutrición
y la excreción sin un intestino ni un riñón funcionales, y las hormonas
que mantienen el embarazo; no puede impedir el paso de todos los tóxicos
--- el alcohol, la nicotina y algunos virus lo atraviesan.
\textbf{20.} Distensión del cuello del útero $\to$ hipotálamo $\to$
oxitocina $\to$ contracción $\to$ más distensión, hasta que el nacimiento
elimina el estímulo. Antes del término, el útero dominado por la
progesterona tiene pocos receptores de oxitocina y ninguna unión
comunicante, de modo que las contracciones no se propagan y el bucle no
puede cerrarse.
\textbf{21.} $750\times 2.9 = \qty{2.2}{MJ}$; coste $2.2/0.8 =
\qty{2.7}{MJ}$.
\textbf{22.} Crecimiento $30\times 6 = \qty{0.18}{MJ}$, mantenimiento
$300\times 4 = \qty{1.2}{MJ}$: \qty{1.4}{MJ}, unas dos terceras partes de
la energía de la leche; el resto se destina a la actividad, al calor y a
las pérdidas.
\textbf{23.} $270\times 2.7 = \qty{730}{MJ}$, dos veces y media el
embarazo.
\textbf{24.} La succión eleva la prolactina y suprime la liberación
pulsátil de la hormona que impulsa la \omterm{def:b2:mammal-reproduction:ovary}{ovulación}, de modo que esta solo
vuelve meses después; sin anticoncepción, esto espacia los nacimientos
dos o tres años.
\textbf{25.} $P = 0.95$; prueba positiva el día 17 tras la fecundación
(día 31 del ciclo); margen de oxígeno de ocho; un día de leche le cuesta a
la madre \qty{2.7}{MJ}.
\end{solution}
